\chapter{What ``kernel-checked'' means}\label{ch:formal}

This chapter explains the mark \Kstatus. It is short, because the idea is
simple, and it is here because the idea is routinely misunderstood in both
directions --- as guaranteeing more than it does, and as guaranteeing less.

\section{The setting}

The statements marked \Kstatus\ in this book live in a Lean~4 development
\cite{deMouraUllrich2021,mathlib4} in the public repository
\artifact{SocrateAI-DualScaleTopologicalUniverseModel-LeanProposal}, at the
tagged release \artifact{v0.25-paper-rankjump} (main branch commit
\texttt{ad64ed2}, 29 September 2026). The lattice interface (Gram matrices,
the hyperbolic plane, $E_8$, reflections, the norm form) is imported from a
second public development~\cite{LeanMaster2026}; nothing in this book depends
on anything else that development contains.

A Lean theorem is a term whose type is the statement; the \emph{kernel} is a
small program that checks the term has that type. When it does, the
statement is a consequence of the axioms the term uses. The developments here
use exactly Lean's three standard axioms (propositional extensionality,
choice, quotient soundness); every theorem cited in this book was checked with
\texttt{\#print axioms} to depend on those alone, with no \texttt{sorry}
(the placeholder that admits a goal) and no \texttt{native\_decide} (which
trusts the compiler rather than the kernel).

\section{What the mark certifies}

Exactly this: \emph{the statement, as written in Lean, follows from the three
axioms.} Not the docstring, not the name, not the paper's paraphrase of it.
Here is one of the statements this book uses, verbatim:

\begin{leanbox}
\texttt{theorem root7\_perp\_iff (x : Fin 3 $\to$ $\Z$) :}\\
\texttt{\quad pairing T7 x root7 = 0 $\leftrightarrow$ $\exists$ a b : $\Z$, x = ![a, 2 * a - 7 * b, b]}

\smallskip\noindent
with \texttt{T7 := TN 7}, \texttt{TN N := !![0, 1, 0; 1, 0, 0; 0, 0, 2 * N]},
\texttt{root7 := ![2, -4, 1]}, and \texttt{pairing G x y := x $\cdot_v$ (G $*_v$ y)}
(the dot product of \texttt{x} with the matrix--vector product \texttt{G y}). The
mathematical symbols are rendered here for the printed page; the source uses the
library's Unicode notation.
\end{leanbox}

Read it against Example~\ref{ex:root7}. It says: $x$ is orthogonal to
$(2,-4,1)$ under the Gram matrix of $\U\oplus\lat{14}$ if and only if $x$ has
the form $(a,2a-7b,b)$. That is the full content of ``$v_2^{\perp}$ is spanned
by $(1,2,0)$ and $(0,-7,1)$'', and it is stronger than ``$(1,2,0)$ and
$(0,-7,1)$ are orthogonal to $v_2$'' --- which would be true of any two vectors
in the complement and would not pin the complement down.

\section{Two ways to certify less than the name says}

The first failure mode is \emph{vacuity}: a statement of the form ``there
exists a polynomial of degree $3$ such that\ldots'' that is true of any cubic.
The development had such statements in an early generation; they were deleted,
and the episode is recorded in its escalation log.

The second is subtler and is the reason this chapter exists: a statement that
is true, compiles, and proves \emph{less than its name}. The recorded
instance: a lemma named for the primitivity of $e+Nf$ in $\U$ whose statement
was $\mathrm{IsCoprime}\,(1:\Z)\,N$ --- true for every $N$ in every commutative
ring, mentioning no vector, no basis and no lattice. Nine such instances were
found in one audit, none by any automated gate; all were found by reading
statements against names. The repaired version states the vector:

\begin{leanbox}
\texttt{theorem glue\_primitive (N : $\Z$) \{d : $\Z$\} \{w : Fin 2 $\to$ $\Z$\}}\\
\texttt{\quad (h : ![1, N] = d $\bullet$ w) : IsUnit d}
\end{leanbox}

The rule the development adopted, and which this book follows in every
\Kstatus\ box, is: \emph{every noun in the name must appear in the statement}.
A reader auditing a formal development should read the statements, not the
names, and not the docstrings.

\section{Producer and verifier}

For the lattice results of Chapter~\ref{ch:lattices} the gates were run twice.
The author's session ran the build, the \texttt{sorry} and axiom audits and a
statement-hash lock, and printed the axioms of each theorem. A second team of
the program, in a different repository, then re-ran the same gates on the exact
commits, recomputed the lattice arithmetic from its own certificate data
before trusting either report, and recorded the result against the commit and
the SHA-256 of the source file. Independently, the development
\cite{LeanMaster2026} stated the same five complements in its own file
(\artifact{DualScaleDyons/RankJump.lean}); the two sets of statements were
compared and agree --- a comparison of statements, not a second verification of
either proof, since neither side built the other's file. We record one detail
because it is the kind of thing a reader should know happens: a control value
$\langle v_2,v_3\rangle$ was first written by hand as $-42$; the kernel
rejected it; the value that compiles is $-154$.

\section{What was not checked}

The identification of $\U\oplus\lat{14}$ with the transcendental lattice of
the $s_7$ family; Nikulin's theorems; Dolgachev's and Doran's theorems; the
Kodaira--Tate table; the Shioda--Tate formula; every $q$-expansion identity;
every numerical recognition. In short: everything marked \Lstatus, \Estatus\ or
\Nstatus\ in this book. The kernel checked $3\times3$ integer matrices and
polynomial identities. That is a great deal less than the theory of K3
surfaces and a great deal more than nothing: it is exactly the part of the
argument where a slip of the pen --- $-42$ for $-154$, $I_1$ for $I_2$, one
rank-3 lattice for another --- would otherwise have gone unnoticed.

\section{Exercises}

\begin{exercise}
Write the statement of Example~\ref{ex:disc3}'s complement in the style of
\artifact{root7\_perp\_iff}, and check that every noun in the name
``\texttt{root7Inf\_perp\_iff}'' occurs in it.
\end{exercise}

\begin{exercise}
Find, in any formal development you have access to, a theorem whose name
promises more than its statement delivers. (This is easier than it sounds.)
\end{exercise}

\section*{Chapter note}

The repository, commit and release are as stated; the two Lean statements are
transcribed verbatim from \artifact{Agora/Geometry/MnLattice.lean} at that
commit. The account of the audit and of the producer/verifier run follows the
repository's lessons-learnt file and the receipt briefs exchanged between the
program's teams, all of which are in the repository.
